\section{Methodology}

\subsection{Research Design}

This study adopts a quantitative research design to examine factors associated with student academic performance. Quantitative research involves examining measurable variables using numerical data and statistical procedures to investigate relationships between variables and test research questions \citep{creswell2023research}.

The study focuses on measurable characteristics of students and uses statistical modelling to investigate the relationship between these characteristics and academic outcomes. This approach is appropriate for the study because student academic performance and the explanatory factors considered can be represented using numerical measures and analysed statistically.

\subsection{Study Variables}

The dependent variable in the study is student academic performance. The explanatory variables represent measurable characteristics that may be associated with academic performance, including study time, attendance, previous academic performance, and access to learning resources. Previous research has identified study behaviour and attendance among factors that may be associated with students' academic outcomes \citep{masnan2025academic,ha2024attendance}.

The variables are considered jointly so that the statistical model can examine the relationship between academic performance and several explanatory factors simultaneously. The dependent variable represents the outcome of interest, while the explanatory variables are included to investigate their individual contributions to variations in academic performance.

For the purpose of the regression analysis, the variables will be represented numerically where appropriate. This allows the observed student-level data to be incorporated into the statistical model and provides a basis for estimating the relationship between academic performance and the selected explanatory variables.

The variables are considered jointly so that the statistical model can examine the relationship between academic performance and several explanatory factors simultaneously.

\subsection{Multiple Linear Regression Model}

Since the study considers several explanatory variables simultaneously, multiple linear regression is used to examine their relationship with student academic performance. This method allows the contribution of each explanatory variable to be estimated while holding the other variables constant \citep{yang2018predicting}.

The general form of the multiple linear regression model is given by:

\begin{equation}
Y_i = \beta_0 + \beta_1 X_{1i} + \beta_2 X_{2i} + \cdots + \beta_p X_{pi} + \varepsilon_i
\end{equation}

where:

\begin{align*}
Y_i & = \text{academic performance of student } i, \\
\beta_0 & = \text{intercept of the regression model}, \\
\beta_j & = \text{coefficient associated with explanatory variable } X_j, \\
X_{ji} & = \text{value of explanatory variable } j \text{ for student } i, \\
\varepsilon_i & = \text{random error term for student } i.
\end{align*}

The coefficient $\beta_j$ represents the expected change in academic performance associated with a one-unit change in the corresponding explanatory variable, holding the other explanatory variables constant.

\subsection{Model Estimation}

The regression coefficients are estimated using the available student-level observations. The fitted model provides estimates of the relationships between the explanatory variables and academic performance and can subsequently be evaluated using appropriate statistical measures \citep{yang2018predicting}.

\subsection{Model Assumptions}

The multiple linear regression model is based on assumptions concerning the relationship between the variables and the error term. These include linearity, independence of observations, constant error variance (homoscedasticity), and normality of the errors \citep{tamhane2020predictive}.

These assumptions will be considered when evaluating the fitted model to determine whether the regression approach provides an appropriate representation of the observed academic performance data.

\subsection{Data Analysis Procedure}

The collected data will be organised and analysed using appropriate statistical procedures. Descriptive statistics will first be used to summarise the characteristics of the study variables. The multiple linear regression model will then be fitted to examine the relationship between student academic performance and the selected explanatory variables.

The estimated regression coefficients will be examined to determine the direction and magnitude of the relationships between the explanatory variables and academic performance. Measures of model fit and relevant statistical tests will also be considered when evaluating the fitted model.

The results obtained from these analyses will provide the basis for assessing the research questions and evaluating the usefulness of the proposed statistical model.